\subsection{Sharpness of the localization radius}
The slack $U-\beta$ can be as large as a sum of $n$ correction terms,
while a single coordinate interval is removed only when its own excess
exceeds the whole slack. The following family shows that the radius
$\Theta(\sqrt{n\kappa}\,h_j)$ of Lemmas~\ref{lem:states}
and~\ref{lem:commonmesh} is therefore attained, so that filtering alone
leaves at least of order $\sqrt{n\kappa}$ nodes per coordinate on uniform
grids; the graded grids reduce this to
$O(\sqrt{\bar\kappa}\log(n+2))$ (Lemma~\ref{lem:states}), which makes
Theorem~\ref{thm:approx} a fixed-parameter bound (Remark~\ref{rem:fpt}).

\begin{proposition}[Sharpness of localization]\label{prop:sharp}
Let $m\ge1$, $n=2m$, $\Lambda>0$ and $0<g\le\Lambda/2$, and on
$X=[-1,1]^n$, with all coordinates continuous and named
$(u_1,v_1,\dots,u_m,v_m)$, let
\[
F(x)=\sum_{k=1}^m\Bigl[\frac\Lambda2(u_k^2+v_k^2)+(\Lambda-2g)u_kv_k\Bigr].
\]
Then $x^*=0$, $\OPT=0$, the smallest valid curvature bounds are
$L_i=\Lambda$, the largest growth constant is $g$, and $\kappa=\Lambda/g$
takes every value in $[2,\infty)$. Let $G$ be a grid for a subbox that
contains $[-R_0,R_0]^n$, and let every $G_i$ contain $0$ with
$w_i(0)\ge h$. Then for every incumbent value $U\ge0$, every coordinate $i$
and every node $a\in G_i$ with
\[
|a|\le\min\Bigl\{R_0,\ h\sqrt{(n-2)\kappa/16}\Bigr\}
\]
we have $m_i(a)\le U$. Hence every grid interval with such an endpoint is
retained.
\end{proposition}

\begin{corollary}[Filtered uniform grids need order $\sqrt{n\kappa}$ nodes]
\label{cor:uniformgrid}
Take the instance of Proposition~\ref{prop:sharp} with $n\ge4$, and run the
stages of a trial with a common mesh (Lemma~\ref{lem:commonmesh}) without a
cap, with $h_j=s2^{-j}=2^{1-j}$, uniform grids (the construction of
Definition~\ref{def:graded} with $\theta=0$) and initial center $0$. Let
$k_0=\lfloor\sqrt{(n-2)\kappa/16}\rfloor$. At every stage the corrected
minimizer is $0$; at every stage $j\ge1$ with $(k_0+1)h_j\le1$ that ends by
filtering, the filtered box $X^{(j+1)}$ contains
$[-(k_0+1)h_j,(k_0+1)h_j]^n$, and the grid of stage $j+1$, if that stage is
run, has at least $4k_0+5\ge\sqrt{(n-2)\kappa}+1$ nodes in every coordinate.
With graded grids of grading $\theta\in(0,1/4]$ instead: if a stage $j$ ends
by filtering, its box $X^{(j)}$ contains $[-R,R]^n$, where
$R=h_j\sqrt{(n-2)\kappa/16}$, and its grids contain $0$ with
$w_i(0)\ge h_j$, then the grid of stage $j+1$, if that stage is run, has at
least $1+\ln(1+2\theta R/h_j)/\ln(1+\theta)$ nodes in every coordinate.
\end{corollary}

Proofs are in Appendix~\ref{app:growth}. Uniform grids also give
accuracy-independent counts without any grading condition: under point
growth, UC (Algorithm~\ref{alg:uc}), which filters unions of uniform cells,
has at most $12(2\sqrt{n\kappa}+1)$ nodes per coordinate at every stage
(Theorem~\ref{thm:cells} with $r=1$ and $\kappa_S=\kappa$). The resulting
running-time bound is polynomial in $I+q$ and $\kappa$ for fixed $p$, but
because of the factor $\sqrt n$ per coordinate it is not a fixed-parameter
bound.
